\section{Preserving the global boundary invariants}\label{sec:boundary}

We verify Invariants~\ref{inv:past} and~\ref{inv:swap}
for a child \(W'=(a',b')\) of \(W=(a,b)\).
All notation is in the working frame of the child edge.
In a same update this is also the parent's tag frame.
In a switch update the old list is not carried; the newly constructed
upsets are used instead.

Preparatory changes to the schedule of \(W\) preserve the parent
invariants. Its boundary triples and all earlier exterior jobs
stay fixed, so Invariant~\ref{inv:past} is unchanged.
Invariant~\ref{inv:swap} already quantifies over every feasible
reordering of \(W\).

\subsection{Excluding distant qualifiers}

\begin{lemma}\label{lem:past}
The updated child list satisfies Invariant~\ref{inv:past}.
\end{lemma}
\begin{proof}
It suffices to prove the assertion before pruning, since pruning only
removes qualifiers. Let \(z\) qualify for the pre-pruning child list and be scheduled
at or before \(a'\).

If \(a<\operatorname{time}(z)\le a'\), the job belongs to the normalized
open side \(E\) of the parent. By \Cref{lem:update-labels}, it is
high at the edge cutoff \(S\). The local separator at \(a'\) puts
it in \(P_{a'}\) when strictly earlier, and membership in the
triple puts it in \(\wP{a'}\) when its time equals \(a'\).

Now suppose \(z\) is at or before \(a\).
The new injection \(D_a\setminus P_b\) cannot contain such a job.
Nor can either switch entry, by \Cref{lem:G} and the strictness
of \(D_a\). Thus the only remaining case is a same-direction
carry qualifier.
The parent's Invariant~\ref{inv:past} supplies a job \(y\in a\)
with \(z=y\) or \(z\prec y\).
The carry is a global upset, so \(y\) also qualifies by carry.
In particular \(a\) is actual. If \(a'=a\), we are done.

Assume \(a'>a\), and write \(a'=s_i\), with edge cutoff \(S=S_i\).
Follow all the strict advances of the walk from \(a\) through
promotion \(i\), including intermediate advances not retained as marks.
Every cutoff used in this portion of the walk is contained in \(S_i\).
On \(W\), the final carry is exactly \(O\setminus S_i\).
Consequently every carry qualifier in \(E\) is high, and has worse
priority than the low exceptions, at every one of these advances.

The first advance needs the parent exchange invariant, because
\(a\) lies outside the normalized open side \(E\).
Consider that first actual advance off \(a\), with earliest exception
\(x\) at its target \(t\).
All \emph{global} predecessors of \(x\) are before \(a\).
To verify this stronger readiness assertion, every predecessor is
temporally before \(t\). If one were at \(a\), then \(x\in D_a\).
If one were strictly between \(a\) and \(t\), it would belong to
the contiguous side \(E\), be low, and descend from \(a\) by
earliestness of the exception. Again \(x\in D_a\), a contradiction.
There is no other possible predecessor at or after \(a\).

Also \(x\in X=W\setminus D_a\), so Invariant~\ref{inv:cones}
puts it among the old qualifiers, outside the starting low set.
It has already been promoted by stage \(i\), so its old key is
strictly below the least unpromoted key at \(S_i\).
By \Cref{lem:carry}, the boundary job \(y\) has old key at least
that least unpromoted key.
The parent's Invariant~\ref{inv:swap}, applied after the current
preparatory reorderings of \(W\), therefore prohibits exchanging
\(x\) and \(y\) in the complete schedule.

If \(y\) had no global successor at or before \(t\), that exchange
would be feasible. All global predecessors of \(x\) lie strictly before \(a\);
moving it earlier preserves its successor constraints, and delaying
\(y\) to \(t\) preserves its predecessor constraints and, under
the supposition, its successor constraints. Thus \(y\) must have
a strict successor by \(t\).
Such a successor is in \(E\) and is still a carry qualifier.

Any carry job strictly between \(a\) and \(t\) must likewise have
a successor by \(t\): otherwise exchange it with \(x\) within
\(E\), contradicting its lexicographic normalization.
The carry is an upset, so these successors remain in the carry.
Iterating along strictly increasing times reaches a carry job
at \(t\). It cannot be \(x\), which is low and lies outside
the final carry.

At every subsequent advance, the local readiness and exchange argument
from \Cref{lem:shared} applies inside \(E\). All carry jobs remain
high at that advance, so each carry job at its starting slot
has a carry descendant at its target. Tracing from the descendant
of \(y\) at the first target eventually gives a descendant of
\(y\) at \(a'\).
Hence \(z\in\wP{a'}\), as required.
\end{proof}

\subsection{Exchanges after future child reorderings}

\begin{lemma}\label{lem:swap}
The pruned child list satisfies Invariant~\ref{inv:swap},
with its quantifier over all feasible reorderings of \(W'\).
\end{lemma}
\begin{proof}
Take any feasible reordering of \(W'\) in its slots.
Suppose \(x\in W'\setminus D_{a'}\), that \(y\in a'\)
qualifies for the child list, and that the new key of \(x\)
is smaller than the new key of \(y\).
Invariant~\ref{inv:cones}, already proved for the child in
\Cref{lem:update-labels}, ensures that \(x\) qualifies.
By \Cref{lem:prune}, the same comparison holds in the new list
before pruning.

First let \(a'>a\). Both \(x\) and \(y\) are in the parent's
normalized side \(E\), and both are high at its edge cutoff \(S\),
by \Cref{lem:update-labels}. The pre-pruning key order agrees with
the priority order used to normalize \(E\). Thus \(x\) has
better priority than \(y\).
The child reordering changes only slots strictly after \(a'\).
A feasible exchange of \(x,y\) in the complete schedule would
therefore give a feasible full schedule of the same \(E\)
with a better prefix at \(a'\), contrary to
\Cref{lem:lex-prefix}.

Now let \(a'=a\).
In the switch case there is no new qualifier at \(a\), so no
such \(y\) exists.
In the same case neither \(x\) nor \(y\) belongs to the new
injection: \(x\notin D_a\), and no job at \(a\) is a strict
descendant of \(a\).
Their pre-pruning comparison therefore uses only carry entries.
By \Cref{lem:carry}, it implies the same old key comparison
in the parent's list.
Moreover \(x\in W\setminus D_a\).
The preparatory reorderings and the arbitrary subsequent reordering
of \(W'\) together form a feasible reordering confined to \(W\).
The parent's Invariant~\ref{inv:swap} rules out the proposed
exchange.
\end{proof}

\begin{proposition}\label{prop:construction-valid}
Starting from any feasible full schedule, the marking and list
construction is defined at every node. Its pruned lists are global
upset lists of length at most five, with fewer than \(40\) leaves
per entry, and satisfy Invariants~\ref{inv:cones}--\ref{inv:swap}.
Every branch terminates.
\end{proposition}
\begin{proof}
Proceed by depth from the root data of \Cref{sec:construction}.
Given valid ancestor and parent data, the center and the two
normalized sides exist by \Cref{cor:any-ideal,lem:reorder}.
The shared walk and \Cref{lem:edge-labels} define the children.
For a same update the new list is direct. For a switch,
\Cref{lem:old-low,lem:G} construct the needed global upset
using only earlier data, as explained at the end of
\Cref{sec:lists}.
Then \Cref{lem:update-labels,lem:list-length,lem:entry-size}
give the first invariant and the bounds.
\Cref{lem:past,lem:swap} give the other two invariants without
using a claim about any deeper child.
This completes the simultaneous induction.
Every child omits the center, so its size decreases strictly.
\end{proof}

Invariant~\ref{inv:past} will let a bounded predicate distinguish
the child from qualifiers arbitrarily far in the past.
Invariant~\ref{inv:swap} is what allows that first invariant to
continue propagating after recursive reorderings. Neither is
replaced by an assertion only about local separators.
